\slajdid{02-s020}
\begin{frame}[label=02-s020]
\gusttekst
Da li su nam uvek na raspolaganju i prave i komplementne vrednosti signala. U većini praktičnih slučajeva NE.

Za formiranje električne, logičke, šeme će nam trebati i invertori pošto nam, na primer, trebaju komplemente vrednosti ulaznih signala jer raspolažemo samo sa pravim vrednostima.
\begin{columns}[c,onlytextwidth]
\begin{column}{.61\textwidth}\centering
\begin{tikzpicture}[DEoznaka,x=10mm,y=6.5mm]
\foreach \y/\low/\lab in {6.5/5.9/C,5.3/4.7/B,4.1/3.5/A}{
 \node[not gate US,draw,minimum width=7mm,minimum height=6mm] (n\lab) at (1.6,\y) {};
 \draw[DEvod] (0,\y) node[left] {\(\lab\)}--(n\lab.input);
 \draw[DEvod] (n\lab.output)--(6.2,\y);
 \draw[DEvod] (.6,\y)--(.6,\low)--(6.2,\low);
}
% Jezgro se ponavlja i na sledećem slajdu sa eksplicitnim invertorima.
\foreach \i/\x in {1/3.1,2/4.4,3/5.7}{
 \node[and gate US,draw,logic gate inputs=nnn,rotate=-90,minimum width=8mm,minimum height=8mm] (g\i) at (\x,2.1) {};
}
\foreach \i/\j/\y in {1/3/6.5,1/2/4.7,1/1/3.5,2/3/5.9,2/2/5.3,2/1/3.5,3/3/5.9,3/2/4.7,3/1/4.1}{
 \draw[DEvod] (g\i.input \j)--(g\i.input \j |- 0,\y);
}
\node[or gate US,draw,logic gate inputs=nnn,rotate=-90,minimum width=8mm,minimum height=8mm] (z) at (4.4,.15) {};
\draw[DEvod] (g1.output)--(g1.output |- 0,.95)-|(z.input 3);
\draw[DEvod] (g2.output)--(z.input 2);
\draw[DEvod] (g3.output)--(g3.output |- 0,.95)-|(z.input 1);
\draw[DEvod] (z.output)--++(0,-.35) node[below] {\(F\)};

% Crvene strelice su napomene za tri opterećenja izvora signala A, nisu vodovi.
\draw[DEcrvena,->] (0,4.45)--(1.1,4.45);
\draw[DEcrvena,->] (0,3.9)--(.4,3.9)--(.4,3.1)--(g1.input 1 |- 0,3.1)--($(g1.input 1)+(0,.15)$);
\draw[DEcrvena,->] (.4,3.1)--(g2.input 1 |- 0,3.1)--($(g2.input 1)+(0,.15)$);
\end{tikzpicture}

\end{column}
\begin{column}{.36\textwidth}
Primer: logičko kolo koje je napravilo promenljivu A „vidi“ tri ulaza u logička kola.

\alert{Sa stanovišta realnih logičkih kola to nije dobro.}
\[
F=\bar CBA+C\bar BA+CB\bar A
\]
\end{column}
\end{columns}
\end{frame}
